\documentclass[journal ]{new-aiaa}
%\documentclass[conf]{new-aiaa} for conference papers
\usepackage[utf8]{inputenc}
\usepackage{textcomp}

\usepackage{graphicx}
\usepackage{amsmath}
\usepackage[version=4]{mhchem}
\usepackage{siunitx}
\usepackage{longtable,tabularx}
\setlength\LTleft{0pt} 

\title{Project 2: Turbojet Engine Design/Off-Design Performance Analysis}

\author{Riley J. Kyle\footnote{Aerospace Engineering Student, MAE}}
\affil{Missouri University of Science and Technology, Rolla, Missouri, 65409, United States}
\author{Dr. David W. Riggins\footnote{Professor of Mechanical and Aerospace Engineering, MAE}}
\affil{Missouri University of Science and Technology, Rolla, Missouri, 65409, United States}


\begin{document}

\maketitle

\begin{abstract}
The purpose of this project is to plot the performance three space of a fixed-area missile-class engine with a simple converging nozzle.
This will be done by varying the fuel flow rate from 20$\% $ to 500$\% $ the on design fuel rate, varying the Mach number from 0 to 2.5.
Additionally, this will be done at sea level, 4000 m, 8000 m, and 12000 m altitude.
\end{abstract}

\section*{Nomenclature}
{\renewcommand\arraystretch{1.0}
\noindent\begin{longtable*}{@{}l @{\quad=\quad} l@{}}
$T$  & temperature $[K]$ \\
$P$ &    pressure $[Pa]$\\
$\rho$ & density $\left[kg/m^3\right]$\\
$T_t$ & total temperature $[K]$\\
$P_t$ & total pressure $[Pa]$\\
$\tau$ & $\frac{T_{t,out}}{T_{t,in}}$\\
$\pi$ & $\frac{P_{t,out}}{P_{t,in}}$\\
$\eta$ & isentropic efficiency\\
$u$ & velocity $\left[m/s \right]$\\
$M$ & mach number\\

\multicolumn{2}{@{}l}{Subscripts}\\
$0$ & freestream\\
$d$ & diffuser\\
$c$ & compressor\\
$b$ & burner\\
$t$ & turbine\\
$n$ & nozzle\\
\end{longtable*}}




\section{Introduction}
\lettrine{T}{he} purpose of this project is to plot the performance three space of a fixed-area missile-class engine with a simple converging nozzle.
This will be done by varying the fuel flow rate from 20$\% $ to 500$\% $ the on design fuel rate, varying the Mach number from 0 to 2.5.
Additionally, this will be done at sea level, 4000 m, 8000 m, and 12000 m altitude.



\section{Completely define the on-design engine fluid dynamics and performance}
\subsection{Methodology}

\subsection{Results}


\section{Develop the complete performance envelope for this specific turbojet engine}
\subsection{Methodology}

\subsection{Results}

\subsubsection{Sea Level}

\subsubsection{4000 m}

\subsubsection{8000 m}

\subsubsection{12000 m}

\section{Conclusion}

\section*{Appendix}



\end{document}
